% =================
% AVISO AO LEITOR — trigger warning formal (não diegético)
% Página oficial antes do início da imersão ficcional
% =================
\cleardoublepage
\thispagestyle{empty}
\begingroup
\setlength{\fboxsep}{8mm}
\begin{center}
{\color{sangue}\scshape\Large AVISO AO LEITOR}\\[1.2em]
{\color{sangue}\rule{0.4\paperwidth}{0.6pt}}\\[2.2em]
\end{center}

\noindent Este livro contém, em sua parte ficcional, representações diretas de:

\begin{itemize}
\itemsep0.4em
\item violência institucional psiquiátrica --- incluindo eletrochoque, lobotomia e cárcere em hospital-colônia (Hospital Colônia de Barbacena, 1903--1994);
\item fome, morte e enterro coletivo de crianças e adultos (secas do Ceará, 1915--1932);
\item descrições médicas e necroscópicas detalhadas, com linguagem clínica explícita;
\item morte, luto, abandono e isolamento de longa duração;
\item narrativa em segunda pessoa que convoca diretamente o leitor a ocupar o lugar do \textit{Paciente 1.263}, com instruções simuladas de leitura.
\end{itemize}

\vspace{1em}

\noindent \textbf{Sobre as instruções simuladas:} nenhuma instrução presente no livro produz efeito físico, químico ou biológico real. A ``contaminação'' é um dispositivo narrativo. As páginas foram impressas com papel e tinta convencionais, sem qualquer substância aplicada. O leitor pode interromper a leitura a qualquer momento, em qualquer fragmento, sem nenhuma consequência além daquela que a própria história sugerir.

\vspace{1em}

\noindent \textbf{Sobre sonhos, sono e sintomas:} não há evidência de que a leitura deste livro induza sonhos específicos, altere o ciclo sono-vigília ou produza qualquer sintoma físico. As descrições de reação corporal são recurso formal, não afirmação empírica.

\vspace{1em}

\noindent \textbf{Sobre os códigos e escalas clínicas:} os protocolos diagnósticos, códigos CID e escalas de avaliação que aparecem na parte ficcional são construção literária. Os transtornos citados são condições clínicas reais, mas sua aplicação a um agente narrativo não tem qualquer respaldo médico. \textbf{Este livro não diagnostica o leitor}: ele simula um diagnóstico, como dispositivo de ficção.

\vspace{1em}

\noindent \textbf{Sobre a distinção ficção/realidade:} os personagens Joaquim Antônio Correia, Dr. Heitor Oliveira, Dra. Lúcia Mendes e a criatura são ficcionais. Os currais do governo (1915--1932) e o Hospital Colônia de Barbacena (1903--1994) são fatos históricos documentados. A distinção completa encontra-se na \textit{Nota Histórica}, ao final do volume (p. \pageref{nota-historica}).

\vspace{1em}

\noindent \textbf{Se você é sensível a algum desses temas:} esta obra pode não ser adequada neste momento. Não há vergonha em adiar a leitura ou abandoná-la. Livros não devem ser provas de resistência.

\vspace{1em}

\noindent \textbf{Se você precisar de apoio:} Centro de Valorização da Vida --- CVV, 188 (ligação gratuita, 24 horas) ou \texttt{cvv.org.br}; Unidades de Saúde Mental (CAPS) em todo o território nacional; ou sua rede de confiança.

\vspace{2.2em}
\begin{center}
{\color{sangue}\rule{0.4\paperwidth}{0.6pt}}\\[1em]
{\small\itshape O Arquivista}
\end{center}
\endgroup
\null\vfill
\clearpage
